\documentclass[10pt,a4paper]{amsart}
\usepackage[margin=19mm]{geometry}
\usepackage{amsmath,amssymb,mathtools}
\usepackage[scr=rsfs,cal=euler]{mathalfa}
\usepackage{xfrac}
\usepackage[unicode,hidelinks,bookmarks=false]{hyperref}
\newcommand{\ie}{i.e.\ }
\newcommand{\cf}{cf.\ }
\newcommand{\resp}{respectively\ }
\newcommand{\Subsection}{\subsection}
\providecommand{\nobreakdash}{\nobreak\hbox{-}}
\setlength{\emergencystretch}{2em}
\multlinegap0pt
\usepackage{amssymb}
\usepackage{xcolor}
\usepackage{algorithm}\usepackage{algpseudocode}
\newtheorem{lemma}{Lemma}
\title{On Minimax Optimality and Uniqueness of Fixed-Step First-Order Methods for Smooth Convex Optimization}
\author{Benjamin Grimmer, Sunghyeon Jo, Chanwoo Park}
\date{Source: arXiv:2609.07671v1 (CC BY 4.0)}
\begin{document}
\maketitle

\noindent\emph{Source and licence.} On Minimax Optimality and Uniqueness of Fixed-Step First-Order Methods for Smooth Convex Optimization. Version arXiv:2609.07671v1 (CC BY 4.0), distributed by arXiv under the Creative Commons Attribution 4.0 International licence, http://creativecommons.org/licenses/by/4.0/, as recorded in the arXiv licence metadata for this exact version and shown on the abstract page https://arxiv.org/abs/2609.07671v1. Full licence notice: \texttt{licenses/arxiv-2609.07671-CC-BY-4.0.txt}.\par\smallskip

\noindent\textit{Editorial scope.} Counted unit: the article's lemma lem:lemni-TC-projection (source0.tex lines 1444--1457). The article's complete original proof of it is delivered with it (source0.tex lines 1459--1678, one \texttt{proof} environment of 220 lines). No mathematics is written, repaired or re-derived: the statement and the proof are the article's own lines, in the article's own notation, and no retained line is substituted or re-flowed.  Retained uncounted as the source-local context the proof needs: lines 281--291; lines 624--629; lines 1198--1202; lines 1292--1308; lines 1334--1338; lines 1360--1373; lines 1418--1423; lines 1426--1437.  Nothing was omitted from any retained span, so the package carries no omission record.  No retained line contains a citation, so the wrapper carries no bibliography and no bibliography entry.  No retained line contains a figure environment, an image-inclusion command or drawing code, so no image is delivered and the package declares no asset dependency.  Editorial formatting only, in addition: an AMS \texttt{amsart} preamble that restates the article's own declarations and macros used by the retained lines, namely the environments \texttt{lemma} on separate counters, together with the standard packages those lines need.  The retained environments print in this document as Lemma 1 in order of appearance, whereas the article prints the same items with the numbering of its own full document.  The complete native source of the article as distributed in the arXiv e-print for this exact version is bundled unchanged as \texttt{sources/209648/source0.tex}.
\begin{lemma}\label{lem:app-inverse-bound}
Let $A$ have nonpositive off-diagonal entries, with
$A\mathbf1\geq0$ and $\mathbf1^\top A\geq0$. Whenever
\[
Au=e_i,
\qquad
Av=e_j
\]
are solvable, every solution has $0\leq u_j\leq\min\{u_i,v_j\}$.
The same statement holds for every principal submatrix of $A$. In addition, $\operatorname{range}(A)=\operatorname{range}(A^\top)$.
\end{lemma}

\begin{equation}\label{eq:ogm-cut-r}
    \sum_{j=0}^i\delta_j
    =T_{i,i}b_i
    +\sum_{k=0}^{i-1}\sum_{j=i+1}^N(-T_{k,j})b_j
    \geq T_{i,i}b_i.
\end{equation}

\begin{equation}\label{eq:lemni-rate-recursion}
    \Omega_N(\rho_i-\rho_{i+1})^2
    =\rho_i(1-\rho_{i+1}^2)
    \qquad \forall 0\leq i\leq N.
\end{equation}

\begin{equation}\label{eq:lemni-TC-problem}
\begin{aligned}
 r_{D\rightarrow G}
 =\inf_{W,T,C,r}\quad &r\\
 \mathrm{s.t.}\quad
 &W\text{ is unit lower triangular},\\
 &T\text{ is upper triangular},\qquad T_{i,i}\geq0,\qquad T_{i,j}\leq0 \quad\forall i<j,\\
 &C\text{ is unit lower triangular},\qquad C_{i,j}\leq0\quad \forall i>j,\\
 &(TC)_{i,j}\leq0\quad \forall i\neq j,\qquad TC\mathbf1\geq0,
 \qquad C^\top T^\top\mathbf1\geq0,\\
 &\begin{pmatrix}
       \frac{\sqrt r}{2} &-\frac12\mathbf1^\top TC\\[1mm]
       -\frac12C^\top T^\top\mathbf1
       &\operatorname{sym}(W^\top TC)-\frac1{2\sqrt r}e_Ne_N^\top
   \end{pmatrix}\succeq0,\qquad r>0.
\end{aligned}
\end{equation}

\begin{equation}\label{eq:lemni-diagonal-conditions}
    (\mathbf1^\top Te_i)^2+(e_N^\top C^{-1}e_i)^2
    \leq 2\sqrt r\,T_{i,i}
    \qquad \forall 0\leq i\leq N.
\end{equation}

\begin{lemma}[The Lemniscate Problem]\label{lem:lemni-scalar-prob}
For every $N\geq1$, the optimization problem
\begin{equation}\label{eq:lemni-scalar-prob}
\begin{aligned}
 \min_{1=\rho_0\geq\rho_1\geq\cdots\geq\rho_{N+1}=0}\quad
 &\max_{0\leq i\leq N}
 \frac{(\rho_i-\rho_{i+1})^2}
 {\rho_i(1-\rho_{i+1}^2)}
\end{aligned}
\end{equation}
has optimal value $1/\Omega_N$, attained uniquely by the sequence in
\eqref{eq:lemni-rate-recursion}. At this sequence, every term in the finite
maximum equals $1/\Omega_N$. 
\end{lemma}

\begin{equation}\label{eq:lemni-C-star}
    (C_L^\star)_{i,i}=1,
    \qquad
    (C_L^\star)_{i,j}=-\frac{\phi_{j+1}-\phi_j}{\phi_i}
    \quad \forall 0\leq j<i\leq N,
\end{equation}

\begin{equation}\label{eq:lemni-T-star}
\begin{aligned}
    (T_L^\star)_{i,i}
    &=\frac{1-\rho_{i+1}^2}{2\rho_{i+1}}
      \frac{\phi_i}{\phi_{i+1}}
      &&\forall 0\leq i<N,\\
    (T_L^\star)_{i,i+1}
    &=-\frac{1-\rho_{i+1}^2}{2\rho_{i+1}}
      &&\forall 0\leq i<N,\\
    (T_L^\star)_{N,N}&=\phi_N.
\end{aligned}
\end{equation}

\begin{lemma}[Projection onto the Lemniscate Problem]\label{lem:lemni-TC-projection}
Let $W,T,C,r$ be feasible in~\eqref{eq:lemni-TC-problem}. Then
\begin{equation}\label{eq:lemni-projection-bound}
    \max_{0\leq i\leq N}
    \frac{(\mathbf1^\top Te_i)^2
    +(e_N^\top C^{-1}e_i)^2}{2T_{i,i}}
    \geq \frac1{\Omega_N}.
\end{equation}
Here the ratio is defined as zero when $T_{i,i}=0$, since
\eqref{eq:lemni-diagonal-conditions} forces its numerator to vanish.
Consequently, $r\geq r_\mathtt{Lemni}$. If $r=r_\mathtt{Lemni}$, then
$T=T_L^\star$, $C=C_L^\star$, and every inequality in
\eqref{eq:lemni-diagonal-conditions} is tight.
\end{lemma}

\begin{proof}
Put
\[
    p=T^\top\mathbf1,
    \qquad
    q=C^{-\top}e_N.
\]
Since $C^{-1}\geq0$ and $C^\top p\geq0$, we have $p\geq0$; also $q\geq0$.
If $T_{i,i}=0$, then the $i$th diagonal condition in
\eqref{eq:lemni-diagonal-conditions} gives $p_i=q_i=0$. As in the preceding
sections, nonnegative column sums imply that the entire $i$th column of $T$ is
zero. The columns indexed by $i$ with $T_{i,i}>0$ are linearly independent, so
\[
    \operatorname{range}(T^\top)
    =\{z:z_i=0\text{ whenever }T_{i,i}=0\}.
\]
Hence $q$ lies in the range of $T^\top$, and $C^\top q=e_N$ shows that
$e_N$ lies in the range of $(TC)^\top$. Lemma~\ref{lem:app-inverse-bound}
gives
\[
\operatorname{range}(TC)=\operatorname{range}((TC)^\top).
\]
Since $C$ is invertible, $e_N$ therefore lies in the range of $T$.
Backward substitution gives a vector $b\geq0$ satisfying $Tb=e_N$, where we
set $b_i=0$ whenever $T_{i,i}=0$.

Here $p,q,b\geq0$, and $q_N=1$. Let the left-hand side
of~\eqref{eq:lemni-projection-bound} be $m$. By
\eqref{eq:lemni-diagonal-conditions}, $m\leq\sqrt r$. Define
\[
    E_{-1}=0,
    \qquad
    E_i=\sum_{j=0}^i p_jb_j
    \quad \forall 0\leq i\leq N,
\]
and
\[
    D_i=\sum_{j=i}^N b_jq_j
    \quad \forall 0\leq i\leq N,
    \qquad
    D_{N+1}=0.
\]
Then
\[
    0=E_{-1}\leq E_0\leq\cdots\leq E_N=1,
    \qquad
    D_0\geq D_1\geq\cdots\geq D_N>D_{N+1}=0.
\]
The same cut expansion as~\eqref{eq:ogm-cut-r} gives
\begin{equation}\label{eq:lemni-forward-cut}
    T_{i,i}b_i\leq E_i
    \qquad \forall 0\leq i\leq N.
\end{equation}

For the reverse inequality, fix $i$ with $T_{i,i}>0$ and consider
the trailing principal block $A_i=(TC)_{i:N,i:N}$. The vectors
\[
u=C_{i:N,i:N}^{-1}b_{i:N},
\qquad
v=\frac1{T_{i,i}}C_{i:N,i:N}^{-1}e_i
\]
solve $A_iu=e_N$ and $A_iv=e_i$. Their relevant coordinates are
$u_i=b_i$, $u_N=D_i$, and $v_i=1/T_{i,i}$.
Applying Lemma~\ref{lem:app-inverse-bound} with the two source indices in the
order $N,i$ gives
$b_i\leq\min\{D_i,1/T_{i,i}\}$. When
$T_{i,i}=0$, our choice gives $b_i=0$, so in all cases
\begin{equation}\label{eq:lemni-reverse-cut}
    b_i\leq D_i
    \qquad \forall 0\leq i\leq N.
\end{equation}

For every $i$, the definition of $m$ gives $p_i^2+q_i^2\leq2T_{i,i}m$.
Multiplying by $b_i^2$ and using
\[
    E_i-E_{i-1}=p_ib_i,
    \qquad
    D_i-D_{i+1}=b_iq_i,
\]
together with~\eqref{eq:lemni-forward-cut}--\eqref{eq:lemni-reverse-cut},
yields
\begin{equation}\label{eq:lemni-path-inequality}
    \frac{(E_i-E_{i-1})^2+(D_i-D_{i+1})^2}
    {2E_iD_i}
    \leq m
    \qquad \forall 0\leq i\leq N.
\end{equation}
If $E_iD_i=0$, the preceding inequality forces the $i$th
step to have zero length (by our convention that $0/0=0$). Thus it suffices to focus only on indices with $E_iD_i>0$.

For $0\leq i\leq N+1$, write $(E_{i-1},D_i)=R_i(\cos\vartheta_i,\sin\vartheta_i)$.
The monotonicity above gives
\[
    \vartheta_0=\frac\pi2
    \geq\vartheta_1\geq\cdots\geq
    \vartheta_{N+1}=0.
\]
The $i$th numerator in~\eqref{eq:lemni-path-inequality} is the squared
distance between the consecutive points with indices $i$ and $i+1$, while its
denominator is
\[
    2E_iD_i
    =2R_iR_{i+1}\cos\vartheta_{i+1}\sin\vartheta_i.
\]
Consequently,
\[
\begin{aligned}
&\frac{(E_i-E_{i-1})^2+(D_i-D_{i+1})^2}{2E_iD_i}\\
&\qquad=
\frac{R_i/R_{i+1}+R_{i+1}/R_i
-2\cos(\vartheta_i-\vartheta_{i+1})}
{2\cos\vartheta_{i+1}\sin\vartheta_i}\\
&\qquad\geq
\frac{1-\cos(\vartheta_i-\vartheta_{i+1})}
{\cos\vartheta_{i+1}\sin\vartheta_i},
\end{aligned}
\]
where the final inequality uses
$R_i/R_{i+1}+R_{i+1}/R_i\geq2$.

Set
\[
    \rho_i=\tan\frac{\vartheta_i}{2}
    \qquad \forall 0\leq i\leq N+1.
\]
Then
\[
    1=\rho_0\geq\rho_1\geq\cdots\geq\rho_{N+1}=0,
\]
and the half-angle identities give
\[
    \frac{1-\cos(\vartheta_i-\vartheta_{i+1})}
    {\cos\vartheta_{i+1}\sin\vartheta_i}
    =
    \frac{(\rho_i-\rho_{i+1})^2}
    {\rho_i(1-\rho_{i+1}^2)}.
\]
Combining this identity with~\eqref{eq:lemni-path-inequality} and
Lemma~\ref{lem:lemni-scalar-prob} proves
\[
\sqrt r\geq m\geq\frac1{\Omega_N},
\]
which is~\eqref{eq:lemni-projection-bound} and gives
$r\geq r_\mathtt{Lemni}$.

Suppose now that $r=r_\mathtt{Lemni}$. The unique scalar optimizer has every
local term equal to $1/\Omega_N$. Since each such term is bounded above by the
corresponding path ratio, which is in turn bounded above by $m\leq\sqrt r$,
every inequality in the preceding chain is tight, and the scalar sequence is
the unique sequence in~\eqref{eq:lemni-rate-recursion}. Equality in the radial inequality gives
$R_i=R_{i+1}$ for every $i$. Since $R_0=D_0$ and
$R_{N+1}=E_N=1$, every radius equals one. The half-angle formulas therefore
give
\begin{equation}\label{eq:lemni-equality-path}
    E_{i-1}=\frac{1-\rho_i^2}{1+\rho_i^2},
    \qquad
    D_i=\frac{2\rho_i}{1+\rho_i^2}=\frac1{\phi_i}.
\end{equation}
Equality in~\eqref{eq:lemni-forward-cut}--\eqref{eq:lemni-reverse-cut} gives
\[
    b_i=D_i=\frac1{\phi_i},
    \qquad
    T_{i,i}b_i=E_i.
\]
In particular every $b_i$ is positive, so $T_{i,i}>0$ for every
$i$ and $T$ is nonsingular. These identities give the diagonal in
\eqref{eq:lemni-T-star}.

Equality in the forward cut forces every entry of $T$ above the first
superdiagonal to vanish. The equation $Tb=e_N$ then gives
\[
    T_{i,i+1}
    =-T_{i,i}\frac{b_i}{b_{i+1}}
    =-\frac{1-\rho_{i+1}^2}{2\rho_{i+1}},
\]
which completes~\eqref{eq:lemni-T-star}.

It remains to determine $C$. Put $y=C\mathbf1$. Since $T$ is a
nonsingular $M$-matrix and $Ty=TC\mathbf1\geq0$, one has
$y\geq0$. The upper bidiagonal form and $Tb=e_N$ show that
\[
    \frac{y_0}{b_0}\geq\frac{y_1}{b_1}\geq\cdots\geq\frac{y_N}{b_N}\geq0.
\]
Here $y_0=b_0=1$. Moreover, $q^\top y=e_N^\top\mathbf1=1$ and $q^\top b=D_0=b_0=1$.
Every $q_i b_i=D_i-D_{i+1}$ is positive, so the preceding monotonicity forces
$y=b$.

For $i>j$, set $\gamma_{i,j}=-C_{i,j}/b_i\geq0$. The identity
$C\mathbf1=b$ and the off-diagonal inequalities $(TC)_{i,j}\leq0$ give
\[
    \sum_{j<i}\gamma_{i,j}=\phi_i-1,
    \qquad
    \gamma_{i+1,j}\leq\gamma_{i,j}
    \quad \forall j<i<N.
\]
We prove by induction on $j$ that
$\gamma_{i,j}=\phi_{j+1}-\phi_j$ for every $i>j$. The row $i=1$ gives
$\gamma_{1,0}=\phi_1-1=\phi_1-\phi_0$. If the assertion is known for columns
below $j$, the row $j+1$ identity gives
$\gamma_{j+1,j}=\phi_{j+1}-\phi_j$, and monotonicity gives
$\gamma_{i,j}\leq\phi_{j+1}-\phi_j$ for all $i>j$. On the other hand,
$C^\top q=e_N$ and $D_i=b_i$ give, with $b_{N+1}=0$,
\[
    q_j=\sum_{i=j+1}^N\gamma_{i,j}(b_i-b_{i+1}),
    \qquad
    q_j=(\phi_{j+1}-\phi_j)b_{j+1},
    \qquad
    \sum_{i=j+1}^N(b_i-b_{i+1})=b_{j+1}.
\]
The strict $\rho$-sequence makes every weight positive, so equality forces
$\gamma_{i,j}=\phi_{j+1}-\phi_j$ for every $i>j$. Thus
\[
    C_{i,j}=-b_i(\phi_{j+1}-\phi_j)
    =-\frac{\phi_{j+1}-\phi_j}{\phi_i},
\]
which is~\eqref{eq:lemni-C-star}. Thus
$T=T_L^\star$ and $C=C_L^\star$. Every local objective
inequality is also tight, so every inequality in
\eqref{eq:lemni-diagonal-conditions} is tight.
\end{proof}

\end{document}
